\section*{Sammendrag}
    I denne oppgaven ble virkemåten til halvlederdioder i ulike kretser undersøkt, derunder klippekretser, enveis og toveis likeretting, samt likeretter med RC-ledd. Hensikten med arbeidet var å oppnå forståelse for hvordan dioder kan benyttes til å forme signaler og omforme vekselspenning til likespenning.

    Forsøket ble utført på elektronikklaboratoriet på Cyberingeniørskolen ved Jørstadmoen leir, og hadde en varighet på omtrent 4 timer. Det ble benyttet oscilloskop, multimeter og signalgenerator for å måle og sammenligne teoretiske og praktiske resultater.

    Resultatene viste at dioder effektivt kan brukes til å klippe bort deler av et signal og til å likerette vekselspenning. Enveis likeretting ga en pulserende likespenning med lav middelverdi, mens toveis likeretting utnyttet begge halvperioder og ga et høyere og jevnere nivå. Videre ble det observert at zenerdioder muliggjør mer nøyaktig spenningsbegrensning i klippekretser. Bruk av kondensator i et RC-ledd førte til en jevnere utgangsspenning, der større kapasitans ga mindre ripple og høyere likespenning.

    Det viste generelt god overensstemmelse mellom teori og måleresultater, men med noen avvik som følge av ikke-ideelle komponenter og måleusikkerhet i instrumentene.
